\documentclass[10pt,a4paper]{amsart}
\usepackage{amsmath,amssymb,amsthm,graphicx}
% Retained minimal preamble for the Stacks Project chapter more-morphisms.tex (More on Morphisms).
% Verbatim excerpt of the chapter preamble: Xy-pic setup (native preamble lines 14-18) and the native
% theorem declarations with the mathematical macros (native preamble lines 154-217). Copied mechanically;
% nothing added, removed or redefined. The wrapper class and the amsmath/amssymb/amsthm/graphicx base
% packages are supplied by the assembler, so the class/externaldocument part of the native preamble is omitted.
% For commutative diagrams we use Xy-pic
\usepackage[all]{xy}

% We use 2cell for 2-commutative diagrams.
\xyoption{2cell}
% Theorem environments.
%
\theoremstyle{plain}
\newtheorem{theorem}[subsection]{Theorem}
\newtheorem{proposition}[subsection]{Proposition}
\newtheorem{lemma}[subsection]{Lemma}

\theoremstyle{definition}
\newtheorem{definition}[subsection]{Definition}
\newtheorem{example}[subsection]{Example}
\newtheorem{exercise}[subsection]{Exercise}
\newtheorem{situation}[subsection]{Situation}

\theoremstyle{remark}
\newtheorem{remark}[subsection]{Remark}
\newtheorem{remarks}[subsection]{Remarks}

\numberwithin{equation}{subsection}

% Macros
%
\def\lim{\mathop{\mathrm{lim}}\nolimits}
\def\colim{\mathop{\mathrm{colim}}\nolimits}
\def\Spec{\mathop{\mathrm{Spec}}}
\def\Hom{\mathop{\mathrm{Hom}}\nolimits}
\def\Ext{\mathop{\mathrm{Ext}}\nolimits}
\def\SheafHom{\mathop{\mathcal{H}\!\mathit{om}}\nolimits}
\def\SheafExt{\mathop{\mathcal{E}\!\mathit{xt}}\nolimits}
\def\Sch{\mathit{Sch}}
\def\Mor{\mathop{\mathrm{Mor}}\nolimits}
\def\Ob{\mathop{\mathrm{Ob}}\nolimits}
\def\Sh{\mathop{\mathit{Sh}}\nolimits}
\def\NL{\mathop{N\!L}\nolimits}
\def\CH{\mathop{\mathrm{CH}}\nolimits}
\def\proetale{{pro\text{-}\acute{e}tale}}
\def\etale{{\acute{e}tale}}
\def\QCoh{\mathit{QCoh}}
\def\Ker{\mathop{\mathrm{Ker}}}
\def\Im{\mathop{\mathrm{Im}}}
\def\Coker{\mathop{\mathrm{Coker}}}
\def\Coim{\mathop{\mathrm{Coim}}}

% Boxtimes
%
\DeclareMathSymbol{\boxtimes}{\mathbin}{AMSa}{"02}

%
% Macros for moduli stacks/spaces
%
\def\QCohstack{\mathcal{QC}\!\mathit{oh}}
\def\Cohstack{\mathcal{C}\!\mathit{oh}}
\def\Spacesstack{\mathcal{S}\!\mathit{paces}}
\def\Quotfunctor{\mathrm{Quot}}
\def\Hilbfunctor{\mathrm{Hilb}}
\def\Curvesstack{\mathcal{C}\!\mathit{urves}}
\def\Polarizedstack{\mathcal{P}\!\mathit{olarized}}
\def\Complexesstack{\mathcal{C}\!\mathit{omplexes}}
% \Pic is the operator that assigns to X its picard group, usage \Pic(X)
% \Picardstack_{X/B} denotes the Picard stack of X over B
% \Picardfunctor_{X/B} denotes the Picard functor of X over B
\def\Pic{\mathop{\mathrm{Pic}}\nolimits}
\def\Picardstack{\mathcal{P}\!\mathit{ic}}
\def\Picardfunctor{\mathrm{Pic}}
\def\Deformationcategory{\mathcal{D}\!\mathit{ef}}
\usepackage{cleveref}
\crefname{originalreference}{original reference}{original references}
\Crefname{originalreference}{Original reference}{Original references}
\makeatletter
\newlabel{coherent-lemma-induction-principle}{{08DR}{source}{}{sourcepointer-1}{}}
\newlabel{coherent-lemma-induction-principle@cref}{{[originalreference][0][]08DR}{[1][0][]source}{}{}{}}
\makeatother
\crefname{originalreference}{original reference}{original references}
\Crefname{originalreference}{Original reference}{Original references}
\makeatletter
\newlabel{morphisms-lemma-characterize-affine}{{01S8}{source}{}{sourcepointer-2}{}}
\newlabel{morphisms-lemma-characterize-affine@cref}{{[originalreference][0][]01S8}{[1][0][]source}{}{}{}}
\makeatother
\crefname{originalreference}{original reference}{original references}
\Crefname{originalreference}{Original reference}{Original references}
\makeatletter
\newlabel{schemes-lemma-characterize-quasi-separated}{{01KO}{source}{}{sourcepointer-3}{}}
\newlabel{schemes-lemma-characterize-quasi-separated@cref}{{[originalreference][0][]01KO}{[1][0][]source}{}{}{}}
\makeatother
\crefname{originalreference}{original reference}{original references}
\Crefname{originalreference}{Original reference}{Original references}
\makeatletter
\newlabel{schemes-lemma-characterize-separated}{{01KP}{source}{}{sourcepointer-4}{}}
\newlabel{schemes-lemma-characterize-separated@cref}{{[originalreference][0][]01KP}{[1][0][]source}{}{}{}}
\makeatother
\title{A scheme has a finite affine stratification exactly when it is quasi-compact and quasi-separated}
\author{The Stacks Project Authors}
\date{stacks-project:more-morphisms.tex:local-snapshot-20261011; GNU FDL 1.2 or later}
\begin{document}
\maketitle
\noindent\textit{Editorial note.} Counted span is the author's own proof environment, source lines 21520-21566, immediately following the statement at 21511-21518.\par\medskip
\noindent\textit{Editorial note.} The proof ends by adding one extra stratum Z' as a minimal element of the partially ordered index set; that construction is displayed inside the span.\par\medskip
\noindent\textit{Editorial note.} Four cross-chapter pointers (Coherent, Morphisms, Schemes) are recorded in assembly.external\_references.\par\medskip
\noindent\textit{Allowed standard prerequisites (editorial).}\begin{itemize}
\item Characterisations of quasi-separated and separated schemes. Schemes, lemmas schemes-lemma-characterize-quasi-separated and schemes-lemma-characterize-separated, cited at proof lines 21524-21530. Turns the intersection behaviour of affine opens into quasi-separatedness, giving the necessity direction.
\item Affineness of a morphism is local on the target. Morphisms, lemma morphisms-lemma-characterize-affine, cited at proof lines 21559-21561. Lets the affineness of the pieces U\_i cap Z and U\_i cap W\_j be checked over the affine opens U and V.
\item Induction principle for quasi-compact quasi-separated schemes. Coherent, lemma coherent-lemma-induction-principle, cited in the proof. Provides the induction step that builds the finite affine stratification for the sufficiency direction.
\end{itemize}\medskip
\ifdefined\hypertarget\hypertarget{sourcepointer-1}{}\fi
\noindent\textit{Source reference.} Original reference 08DR: focused-seeds/geometry/S9-morph-qc-affine-stratification/ref-tags.txt, source lines 4--4. Navigation pointer only: the counted proof cites this statement by \textbackslash{}ref as a standing input; the deduction the counted proof has to make is carried out inside the counted span, and no part of that deduction is delegated to this pointer.\par\medskip
\ifdefined\hypertarget\hypertarget{sourcepointer-2}{}\fi
\noindent\textit{Source reference.} Original reference 01S8: focused-seeds/geometry/S9-morph-qc-affine-stratification/ref-tags.txt, source lines 5--5. Navigation pointer only: the counted proof cites this statement by \textbackslash{}ref as a standing input; the deduction the counted proof has to make is carried out inside the counted span, and no part of that deduction is delegated to this pointer.\par\medskip
\ifdefined\hypertarget\hypertarget{sourcepointer-3}{}\fi
\noindent\textit{Source reference.} Original reference 01KO: focused-seeds/geometry/S9-morph-qc-affine-stratification/ref-tags.txt, source lines 6--6. Navigation pointer only: the counted proof cites this statement by \textbackslash{}ref as a standing input; the deduction the counted proof has to make is carried out inside the counted span, and no part of that deduction is delegated to this pointer.\par\medskip
\ifdefined\hypertarget\hypertarget{sourcepointer-4}{}\fi
\noindent\textit{Source reference.} Original reference 01KP: focused-seeds/geometry/S9-morph-qc-affine-stratification/ref-tags.txt, source lines 7--7. Navigation pointer only: the counted proof cites this statement by \textbackslash{}ref as a standing input; the deduction the counted proof has to make is carried out inside the counted span, and no part of that deduction is delegated to this pointer.\par\medskip
\begin{lemma}
\label{lemma-qc-affine-stratification}
Let $X$ be a scheme. The following are equivalent
\begin{enumerate}
\item $X$ has a finite affine stratification, and
\item $X$ is quasi-compact and quasi-separated.
\end{enumerate}
\end{lemma}

\begin{proof}
Let $X = \bigcup X_i$ be a finite affine stratification. Since each $X_i$
is affine hence quasi-compact, we conclude that $X$ is quasi-compact.
Let $U, V \subset X$ be affine open. Then $U \cap X_i$ and $V \cap X_i$
are affine open in $X_i$ since $X_i \to X$ is an affine morphism. Hence
$U \cap V \cap X_i$ is an affine open of the affine scheme $X_i$ (see
Schemes, Lemma \ref{schemes-lemma-characterize-separated} for example).
Therefore $U \cap V = \coprod U \cap V \cap X_i$ is quasi-compact
as a finite union of affine strata. We conclude that $X$ is
quasi-separated by Schemes, Lemma
\ref{schemes-lemma-characterize-quasi-separated}.

\medskip\noindent
Assume $X$ is quasi-compact and quasi-separated. We may use the
induction principle of
Cohomology of Schemes, Lemma \ref{coherent-lemma-induction-principle}
to prove the assertion that $X$ has a finite affine stratification.
If $X$ is empty, then it has an empty affine stratification.
If $X$ is nonempty affine then it has an affine stratification
with one stratum. Next, assume $X = U \cup V$ where $U$ is quasi-compact open,
$V$ is affine open, and we have a finite affine stratifications
$U = \bigcup_{i \in I} U_i$ and $U \cap V = \coprod_{j \in J} W_j$.
Denote $Z = X \setminus V$ and $Z' = X \setminus U$.
Note that $Z$ is closed in $U$ and $Z'$ is closed in $V$.
Observe that $U_i \cap Z$ and $U_i \cap W_j = U_i \times_U W_j$ are affine
schemes affine over $U$. (Hints: use that $U_i \times_U W_j \to W_j$ is
affine as a base change of $U_i \to U$, hence $U_i \cap W_j$ is affine,
hence $U_i \cap W_j \to U_i$ is affine, hence $U_i \cap W_j \to U$ is affine.)
It follows that
$$
U = \coprod\nolimits_{i \in I} (U_i \cap Z) \amalg
\coprod\nolimits_{(i, j) \in I \times J} (U_i \cap W_j)
$$
is a finite affine stratification with partial ordering on
$I \amalg I \times J$ given by $i' \leq (i, j) \Leftrightarrow i' \leq i$
and $(i', j') \leq (i, j) \Leftrightarrow i' \leq i$ and $j' \leq j$.
Observe that $(U_i \cap Z) \times_X V = \emptyset$ and
$(U_i \cap W_j) \times_X V = U_i \cap W_j$ are affine.
Hence the morphisms $U_i \cap Z \to X$ and $U_i \cap W_j \to X$
are affine because we can check affineness of a morphism
locally on the target
(Morphisms, Lemma \ref{morphisms-lemma-characterize-affine})
and we have affineness over both $U$ and $V$.
To finish the proof we take the stratification above
and we add one additional stratum, namely $Z'$,
whose index we add as a minimal element to the partially ordered set.
\end{proof}

\end{document}
